% !TeX root = ../../Book4.tex
% 纲要标题：## 第5章　幺半范畴、闭结构与对偶
% 纲要标签：b4:ch:monoidal

\chapter{幺半范畴、闭结构与对偶}
\label{b4:ch:monoidal}
\BookChapterNumber{b4:ch:monoidal}

\begin{chapterabstract}
本章以张量积的结合与单位同构建立幺半范畴，说明闭结构如何将带参数的映射表示为内部 Hom，研究求值与余求值的蛇形恒等式，并证明交换环模可对偶恰好等价于有限生成投射性。最后用代数对象和富集范畴统一乘法、模作用与带加法的态射复合。
\end{chapterabstract}

\begin{learninggoals}
\begin{itemize}
\item 写出张量的结合、单位与交换约束，核验五边形、三角形和六边形等式的两条路径。
\item 用求值和柯里化建立内部 Hom 的泛性质，区分内部 Hom 对象与外部态射集合。
\item 构造有限生成投射模的求值与余求值，解释蛇形恒等式怎样产生有限对偶基。
\item 将普通代数、模作用和单子写成相应幺半范畴中的代数结构，辨认预加性范畴的富集数据。
\end{itemize}
\end{learninggoals}

给两个向量空间作张量积，可以把双线性映射变成线性映射。加入第三个空间以后，括号有两种摆法；若再将各空间换成同构的空间，改括号与施行同构又有两种先后次序。此前的泛性质保证这些操作相容。把参与其中的同构和等式单独写出来，就能看出：不只有向量空间，集合的积、模的张量积、函子的复合，都具有类似的结合与单位问题。

\ProseParagraphBreak

另一个问题来自线性映射本身。固定一个空间 \(X\)，给出 \(A\otimes X\to Y\) 的映射，等价于对每个参数 \(a\) 给一个 \(X\to Y\) 的映射。这些映射能否作为同一范畴中的对象继续参与运算？有限维空间还有更特殊的现象：恒等映射可以由有限个向量与泛函的配对恢复，因此 Hom 空间可以用对偶与张量直接表示。无限维空间和有挠元的模会使其中某些步骤失败。闭结构与可对偶性分别回答这两个问题，所需条件并不相同。

\ProseParagraphBreak

本章使用前三章的函子、自然变换、泛性质与伴随，以第二卷第3.1节的自由模、第6.1--6.2节的张量积及第7.1节的投射模为具体基础。结合与单位约束的定义采用 Etingof 等人的《Tensor Categories》中的传统版本\cite[Definition 2.2.8, p. 25]{EtingofEtAl2015TensorCategories}；对偶及其伴随公式可对照\cite[Section 2.10, pp. 40--43, especially Proposition 2.10.8]{EtingofEtAl2015TensorCategories}。该书还研究带额外有限性与正合性条件的张量范畴，本章仅在所需位置提出这些条件。

% !TeX root = ../../Book4.tex
% 纲要标题：### 5.1 幺半结构与张量的相容性
% 纲要标签：b4:sec:monoidal-structure
% 纲要细目（写作范围）：
% #### 5.1.1 结合约束与单位约束
% #### 5.1.2 模、集合与幺半函子
% #### 5.1.3 辫状结构与对称结构

\section{幺半结构与张量的相容性}
\label{b4:sec:monoidal-structure}

三个模的张量积可以写成 \((M\otimes N)\otimes P\)，也可以写成 \(M\otimes(N\otimes P)\)。两者的元素构造不同，却有把 \((m\otimes n)\otimes p\) 送到 \(m\otimes(n\otimes p)\) 的典范同构。四个模出现时，改括号已有数条途径；若它们给出不同映射，省略括号就会隐藏真正的歧义。幺半结构把这些改括号与插入单位的映射列为数据，并规定它们必须满足的相容等式。

\subsection{结合约束与单位约束}
\label{b4:subsec:monoidal-constraints}

\begin{definition}\label{b4:def:monoidal-category}
设 \(\mathcal C\) 为范畴，给定双函子 \(\otimes:\mathcal C\times\mathcal C\to\mathcal C\)、对象 \(I\)，以及关于所列对象自然的同构
\[
\alpha_{X,Y,Z}:(X\otimes Y)\otimes Z\longrightarrow X\otimes(Y\otimes Z),
\qquad
\lambda_X:I\otimes X\longrightarrow X,
\qquad
\rho_X:X\otimes I\longrightarrow X.
\]
若对任意对象 \(W,X,Y,Z\)，有五边形等式
\begin{equation}\label{b4:eq:monoidal-pentagon}
\begin{aligned}
\alpha_{W,X,Y\otimes Z}\circ\alpha_{W\otimes X,Y,Z}
={}&(1_W\otimes\alpha_{X,Y,Z})\circ\alpha_{W,X\otimes Y,Z}
\circ(\alpha_{W,X,Y}\otimes1_Z),
\end{aligned}
\end{equation}
且对任意 \(X,Y\)，有三角等式
\begin{equation}\label{b4:eq:monoidal-triangle}
(1_X\otimes\lambda_Y)\circ\alpha_{X,I,Y}
=\rho_X\otimes1_Y,
\end{equation}
则称这组数据为一个\term[yaoban-fanchou]{幺半范畴}[monoidal category]。\(I\) 称为单位对象，\(\alpha\) 称为结合约束，\(\lambda,\rho\) 称为单位约束。若这些约束都是恒等态射，就称幺半范畴为严格的。
\end{definition}

五边形等式的两边都从 \(((W\otimes X)\otimes Y)\otimes Z\) 到 \(W\otimes(X\otimes(Y\otimes Z))\)：左边改括号两次，右边改三次。它规定的是映射相等，并非只说两端对象同构。三角等式比较两种删除中间单位的方法：先改括号再删 \(I\)，或直接删 \(I\)。下面的图显示后一要求。
\[
\begin{tikzcd}[column sep=large]
(X\otimes I)\otimes Y \ar[rr,"\alpha_{X,I,Y}"] \ar[dr,"\rho_X\otimes1_Y"']
&&X\otimes(I\otimes Y)\ar[dl,"1_X\otimes\lambda_Y"]\\
&X\otimes Y&
\end{tikzcd}
\]
自然性还有另一项职责。例如，对 \(f:X\to X'\)、\(g:Y\to Y'\)、\(h:Z\to Z'\)，它要求
\[
\alpha_{X',Y',Z'}\circ((f\otimes g)\otimes h)
=(f\otimes(g\otimes h))\circ\alpha_{X,Y,Z}.
\]
因此先施行同态再改括号，与先改括号再施行同态，结果相同。

\subsection{模、集合与幺半函子}
\label{b4:subsec:monoidal-examples-functors}

第二卷第6.1--6.2节已经构造张量积及其结合、单位同构。把这些同构放在一起，可以直接检验上述公理。

\begin{proposition}\label{b4:prop:monoidal-modules-sets}
设 \(R\) 为交换含幺环。左 \(R\) 模范畴在 \(\otimes_R\) 下为幺半范畴，单位对象为 \(R\)，其约束由
\[
\begin{aligned}
\alpha((m\otimes n)\otimes p)&=m\otimes(n\otimes p),\\
\lambda(r\otimes m)&=rm,\qquad \rho(m\otimes r)=rm
\end{aligned}
\]
给出。集合范畴在笛卡尔积 \(\times\) 下为幺半范畴，单位对象为单点集，约束是相应的有序对重排。
\end{proposition}
\begin{proof}
对模，交换性使每个左 \(R\) 模同时成为右 \(R\) 模，规定 \(mr=rm\)。纯张量公式保持各处平衡关系。例如
\((rm\otimes n)\otimes p\) 与 \((m\otimes rn)\otimes p\) 的像相同；第二处平衡关系也由 \(m\otimes(rn\otimes p)=m\otimes(n\otimes rp)\) 保持。由张量积泛性质得到 \(\alpha\)，反向重括号公式给出其逆。单位同构的逆分别为 \(m\mapsto1\otimes m\)、\(m\mapsto m\otimes1\)。这些公式均为 \(R\) 线性，并与模同态相容，故约束自然。

五边形的两条路径都把 \(((w\otimes x)\otimes y)\otimes z\) 送到 \(w\otimes(x\otimes(y\otimes z))\)。三角形的两条路径分别把 \((m\otimes r)\otimes n\) 送到 \(m\otimes rn\) 与 \(rm\otimes n\)，由平衡关系相等。纯张量生成各张量积，故相等不只在所列元素上成立，而是态射相等。

对集合，约束及逆由 \(((x,y),z)\leftrightarrow(x,(y,z))\) 和 \((*,x)\leftrightarrow x\leftrightarrow(x,*)\) 给出。四元组的两种重括号得到同一四个坐标，删除单点集也不改变其余坐标，故五边形、三角形及自然性逐坐标成立。
\end{proof}

模范畴的单位不能擅自换成底集合的单点集；它必须是模对象，而且 \(R\otimes_RM\cong M\)。在交换群范畴中，张量的单位因而是 \(\mathbb Z\)，并不是零群。另一方面，直接和 \(\oplus\) 也给出幺半结构，却以零模为单位。这两种幺半结构包含不同的数据。

\ProseParagraphBreak

若 \(R\) 不交换，一般左 \(R\) 模没有与左作用相容的典范右 \(R\) 作用，因而 \(M\otimes_RN\) 不能对任意两个左模直接定义。此时可以改用 \(R\)-\(R\) 双模范畴：张量积仍是 \(R\)-\(R\) 双模，单位为 \(R\)。结合与单位公式的核验相同，但通常没有交换两个因子的典范同构。

\begin{definition}\label{b4:def:strong-monoidal-functor}
设 \(\mathcal C,\mathcal D\) 为幺半范畴，单位分别为 \(I_{\mathcal C},I_{\mathcal D}\)。给定函子 \(F:\mathcal C\to\mathcal D\)、自然同构
\[
F_{2;X,Y}:FX\otimes FY\longrightarrow F(X\otimes Y),
\qquad F_0:I_{\mathcal D}\longrightarrow F(I_{\mathcal C}).
\]
若满足
\[
\begin{gathered}
F(\alpha_{X,Y,Z})\circ F_{2;X\otimes Y,Z}\circ(F_{2;X,Y}\otimes1_{FZ})\\
=F_{2;X,Y\otimes Z}\circ(1_{FX}\otimes F_{2;Y,Z})
\circ\alpha_{FX,FY,FZ},\\
F(\lambda_X)\circ F_{2;I_{\mathcal C},X}\circ(F_0\otimes1_{FX})
=\lambda_{FX},\\
F(\rho_X)\circ F_{2;X,I_{\mathcal C}}\circ(1_{FX}\otimes F_0)
=\rho_{FX},
\end{gathered}
\]
则称 \((F,F_2,F_0)\) 为\term[qiang-yaoban-hanzi]{强幺半函子}[strong monoidal functor]。
\end{definition}

第一式要求 \(F\) 尊重改括号的方式；后两式要求它尊重单位。这里只要求函子带有这些比较态射，不能从 \(FX\otimes FY\) 与 \(F(X\otimes Y)\) 逐对象同构推断它具有幺半结构。上述定义采用的结合与单位约束可对照\cite[Definition 2.2.8, p. 25; Section 2.4, pp. 30--32]{EtingofEtAl2015TensorCategories}。

\subsection{辫状结构与对称结构}
\label{b4:subsec:monoidal-braiding}

改括号保留因子的次序。交换两个因子是另一项结构；对某些代数的表示，普通交换映射甚至可能不保持代数的作用。

\begin{definition}\label{b4:def:braided-symmetric-category}
设 \(\mathcal C\) 为幺半范畴。若自然同构 \(c_{X,Y}:X\otimes Y\to Y\otimes X\) 对任意 \(X,Y,Z\) 满足
\begin{equation}\label{b4:eq:monoidal-hexagons}
\begin{aligned}
c_{X,Y\otimes Z}
={}&\alpha_{Y,Z,X}^{-1}\circ(1_Y\otimes c_{X,Z})\circ\alpha_{Y,X,Z}
\\
&\quad\circ(c_{X,Y}\otimes1_Z)\circ\alpha_{X,Y,Z}^{-1},\\
c_{X\otimes Y,Z}
={}&\alpha_{Z,X,Y}\circ(c_{X,Z}\otimes1_Y)\circ\alpha_{X,Z,Y}^{-1}
\\
&\quad\circ(1_X\otimes c_{Y,Z})\circ\alpha_{X,Y,Z},
\end{aligned}
\end{equation}
则称 \(c\) 为辫状结构，\((\mathcal C,c)\) 为\term[bianzhuang-yaoban-fanchou]{辫状幺半范畴}[braided monoidal category]。这两式称为六边形等式。若还满足 \(c_{Y,X}\circ c_{X,Y}=1_{X\otimes Y}\)，则称其为\term[duichen-yaoban-fanchou]{对称幺半范畴}[symmetric monoidal category]。
\end{definition}

第一条六边形等式说：把 \(X\) 跨过整体 \(Y\otimes Z\)，等于依次跨过 \(Y\) 与 \(Z\)；第二条说明如何把整体 \(X\otimes Y\) 跨过 \(Z\)。公式中的结合约束使每一步都有确定的定义域和陪域。这一版本见\cite[Definition 8.1.1, p. 195]{EtingofEtAl2015TensorCategories}。

\begin{example}\label{b4:exm:symmetric-module-tensor}
交换含幺环 \(R\) 的模范畴取 \(c(m\otimes n)=n\otimes m\)。因为 \(rn\otimes m=n\otimes rm\)，该公式保持平衡关系。它自然，且交换两次恢复原纯张量。第一条六边形等式的两边都把 \(m\otimes(n\otimes p)\) 送到 \((n\otimes p)\otimes m\)；第二条的两边都把 \((m\otimes n)\otimes p\) 送到 \(p\otimes(m\otimes n)\)。纯张量生成各定义域，故两条等式均作为态射等式成立。这给出对称结构。集合的笛卡尔积用交换坐标给出同样的例子。
\end{example}

\begin{example}\label{b4:exm:graded-braiding}
设 \(k\) 为域，\(q\in k^\times\)。在 \(\mathbb Z\) 分次向量空间范畴中，令
\[
(V\otimes W)_d=\bigoplus_{m+n=d}V_m\otimes_kW_n,
\qquad c(v_m\otimes w_n)=q^{mn}w_n\otimes v_m,
\]
态射取保持次数的线性映射，单位 \(k\) 集中在零次。该映射的逆在 \(W_n\otimes V_m\) 上乘 \(q^{-mn}\)。它自然；两条六边形分别化为
\(q^{m(n+r)}=q^{mn}q^{mr}\)、\(q^{(m+n)r}=q^{mr}q^{nr}\)，故给出辫状结构。交换两次乘 \(q^{2mn}\)，所以这一结构对所有分次向量空间对称，当且仅当 \(q^2=1\)。特别地，\(q=-1\) 给出符号 \((-1)^{mn}\)。这里尚没有微分；符号先由次数本身决定。
\end{example}

% !TeX root = ../../Book4.tex
% 纲要标题：### 5.2 闭结构与内部 Hom
% 纲要标签：b4:sec:monoidal-closed
% 纲要细目（写作范围）：
% #### 5.2.1 求值与柯里化
% #### 5.2.2 集合的指数对象与模的内部 Hom
% #### 5.2.3 态射集合与内部对象

\section{闭结构与内部 Hom}
\label{b4:sec:monoidal-closed}

一个双线性映射 \(b:M\times X\to Y\)，既可看成张量积上的线性映射，也可把 \(m\) 固定，得到一个 \(X\to Y\) 的线性映射。于是 \(m\) 对应的并非 \(Y\) 中单个元素，而是一个同态。要把这个对应再写成范畴内部的态射，全部 \(X\to Y\) 的同态必须组织成适当的对象。闭结构正是对这一要求的抽象。

\subsection{求值与柯里化}
\label{b4:subsec:closed-evaluation}

\begin{definition}\label{b4:def:closed-monoidal-category}
设 \((\mathcal C,\otimes,I)\) 为局部小幺半范畴。若对每个对象 \(X\)，函子 \(-\otimes X:\mathcal C\to\mathcal C\) 有右伴随，则称 \(\mathcal C\) 在这一侧为\term[bi-yaoban-fanchou]{闭幺半范畴}[closed monoidal category]。其右伴随在 \(Y\) 处的值称为从 \(X\) 到 \(Y\) 的\term[neibu-Hom]{内部 Hom}[internal Hom]，记为 \(\underline{\operatorname{Hom}}(X,Y)\)，本节也简记为 \([X,Y]\)。具体地，有关于 \(A,Y\) 自然的双射
\begin{equation}\label{b4:eq:closed-adjunction}
\operatorname{Hom}_{\mathcal C}(A\otimes X,Y)
\cong\operatorname{Hom}_{\mathcal C}(A,[X,Y]).
\end{equation}
在对称幺半范畴中，交换两个因子的同构将左右两侧的闭结构相互转换，此时简称闭幺半范畴。
\end{definition}

这里 \([X,Y]\) 是 \(\mathcal C\) 中的对象，等式两侧的 Hom 才是集合。定义并不要求 \([X,Y]\) 的底集合恰好是全体态射；一般范畴甚至没有预先指定的底集合函子。

\begin{proposition}\label{b4:prop:closed-evaluation-transpose}
设 \(\mathcal C\) 为局部小闭幺半范畴。将 \(1_{[X,Y]}\) 经\eqref{b4:eq:closed-adjunction}的逆送回，得到\term[qiuzhi-yingshe]{求值映射}[evaluation map]
\[
\operatorname{ev}_{X,Y}:[X,Y]\otimes X\longrightarrow Y.
\]
对任意 \(f:A\otimes X\to Y\)，存在唯一态射 \(\widehat f:A\to[X,Y]\)，使
\begin{equation}\label{b4:eq:closed-transpose}
f=\operatorname{ev}_{X,Y}\circ(\widehat f\otimes1_X).
\end{equation}
该态射称为 \(f\) 的转置，也称柯里化。内部 Hom 可以组织为双函子
\[
\underline{\operatorname{Hom}}:\mathcal C^{\mathrm{op}}\times\mathcal C\longrightarrow\mathcal C.
\]
对 \(u:X'\to X\)、\(v:Y\to Y'\)，其态射 \([u,v]:[X,Y]\to[X',Y']\) 唯一满足
\begin{equation}\label{b4:eq:closed-hom-functor}
\operatorname{ev}_{X',Y'}\circ([u,v]\otimes1_{X'})
=v\circ\operatorname{ev}_{X,Y}\circ(1_{[X,Y]}\otimes u).
\end{equation}
\end{proposition}
\begin{proof}
记伴随双射为 \(\Phi\)。因 \(\Phi(\operatorname{ev}_{X,Y})=1_{[X,Y]}\)，关于 \(A\) 的自然性给出
\[
\Phi\bigl(\operatorname{ev}_{X,Y}\circ(a\otimes1_X)\bigr)=a
\]
对每个 \(a:A\to[X,Y]\) 成立。因此取 \(a=\Phi(f)\) 得到所需等式，而双射性给出唯一性。

对 \(u,v\)，将\eqref{b4:eq:closed-hom-functor}右边作为从 \([X,Y]\otimes X'\) 到 \(Y'\) 的态射，取其转置，就得到 \([u,v]\)。若 \(u,v\) 都是恒等态射，求值等式由 \(1_{[X,Y]}\) 满足，故 \([1_X,1_Y]=1_{[X,Y]}\)。再给 \(u':X''\to X'\)、\(v':Y'\to Y''\)，连续应用两次求值等式，得到
\[
\begin{aligned}
\operatorname{ev}_{X'',Y''}\circ\bigl(([u',v']\circ[u,v])\otimes1_{X''}\bigr)
=v'v\circ\operatorname{ev}_{X,Y}\circ(1_{[X,Y]}\otimes uu').
\end{aligned}
\]
右边也是 \([uu',v'v]\) 的定义等式。因此转置的唯一性说明
\([u',v']\circ[u,v]=[uu',v'v]\)，证明双函子性。固定 \(X\) 时，该作用正是给定右伴随在态射上的作用，因为伴随在目标变量自然。
\end{proof}

式\eqref{b4:eq:closed-transpose}把“对每个参数给一个映射”变成一个唯一态射。它也是内部 Hom 的泛性质：若另一个对象带求值映射并满足同一存在唯一性，则由\cref{b4:prop:representing-unique}得到保持求值的唯一同构。

\subsection{集合的指数对象与模的内部 Hom}
\label{b4:subsec:closed-sets-modules}

有限积本身具有典范的结合与单位约束。结合约束保持三个投影，单位约束删除终对象的分量；五边形中的两条路径保持同样四个投影，故由积的泛性质相等，三角形也如此。因此有有限积的范畴自然带有以积为张量的幺半结构。

\begin{definition}\label{b4:def:cartesian-closed}
设 \(\mathcal C\) 为局部小范畴，具有有限积。若以积 \(\times\) 及终对象为单位的幺半结构是闭的，即每个 \(-\times X\) 有右伴随，则称 \(\mathcal C\) 为\term[dikaer-bi-fanchou]{笛卡尔闭范畴}[cartesian closed category]。其内部 Hom 通常记为 \(Y^X\)，称为指数对象。
\end{definition}

\begin{example}\label{b4:exm:sets-cartesian-closed}
集合范畴中，令 \(Y^X\) 为全部函数 \(X\to Y\) 的集合，求值为
\[
Y^X\times X\longrightarrow Y,\qquad(h,x)\longmapsto h(x).
\]
给定函数 \(f:A\times X\to Y\)，令 \(\widehat f(a)\) 为函数 \(x\mapsto f(a,x)\)。它满足求值等式；反过来，该等式强制 \(\widehat f(a)(x)=f(a,x)\)，所以唯一。对参数映射 \(a':A'\to A\)、\(u:X'\to X\)、\(v:Y\to Y'\)，对应公式变为
\[
a\longmapsto\bigl(x'\longmapsto v(f(a'(a),u(x')))\bigr),
\]
无论先改变参数再柯里化，还是先柯里化再前后复合，都得到这同一个函数。因此伴随自然，集合范畴笛卡尔闭。
\end{example}

\begin{theorem}[模的 Hom--张量闭结构]\label{b4:thm:closed-module-hom-tensor}
设 \(R\) 为交换含幺环，\(X,Y\) 为左 \(R\) 模。赋予 \(\operatorname{Hom}_R(X,Y)\) 点态的加法与数乘，则它是 \(R\) 模，并且与求值
\[
\operatorname{Hom}_R(X,Y)\otimes_RX\longrightarrow Y,
\qquad h\otimes x\longmapsto h(x)
\]
一起构成 \(R\)-模张量结构的内部 Hom。具体地，对任意左 \(R\) 模 \(A\)，有自然双射
\begin{equation}\label{b4:eq:closed-module-currying}
\begin{gathered}
\operatorname{Hom}_R(A\otimes_RX,Y)
\cong\operatorname{Hom}_R\bigl(A,\operatorname{Hom}_R(X,Y)\bigr),\\
f\longmapsto\bigl(a\longmapsto(x\longmapsto f(a\otimes x))\bigr).
\end{gathered}
\end{equation}
\end{theorem}
\begin{proof}
若 \(h\) 为 \(R\) 线性映射，则 \((rh)(sx)=rsh(x)=s(rh)(x)\)，这里用了 \(R\) 交换，故点态数乘仍给出线性映射。求值对两个变量双线性，并满足
\((rh)(x)=r h(x)=h(rx)\)，所以经张量积得到 \(R\) 线性求值态射。

给定 \(f:A\otimes_RX\to Y\)，对固定 \(a\)，映射 \(x\mapsto f(a\otimes x)\) 线性；而 \(a\mapsto\widehat f(a)\) 也线性，因为
\[
\widehat f(ra)(x)=f(ra\otimes x)=r f(a\otimes x)=\bigl(r\widehat f(a)\bigr)(x).
\]
反过来，给定线性 \(g:A\to\operatorname{Hom}_R(X,Y)\)，映射 \((a,x)\mapsto g(a)(x)\) 双线性，且
\(g(ra)(x)=r g(a)(x)=g(a)(rx)\)。它保持平衡关系，故唯一诱导 \(f_g:A\otimes_RX\to Y\)。两种构造的复合在每个 \(a,x\) 上恢复原值；纯张量生成张量积，所以它们互逆。

对 \(a':A'\to A\)、\(u:X'\to X\)、\(v:Y\to Y'\)，同样直接得到
\[
\widehat{\,v\circ f\circ(a'\otimes u)\,}(a)(x')
=v\bigl(\widehat f(a'(a))(u(x'))\bigr).
\]
这验证三变量的相容性，特别给出伴随所需的两变量自然性。
\end{proof}

该定理不要求 \(X\) 有限生成、投射或 \(Y\) 内射。存在内部 Hom 与 Hom 函子正合是不同的问题。它也不把模范畴变成笛卡尔闭范畴：这里参与伴随的是 \(\otimes_R\)，而模的有限积是直接和。一般非交换环上的 Hom--张量伴随仍可在适当的双模之间成立，但不能因此说全部左模组成其自身的张量闭范畴。

\subsection{态射集合与内部对象}
\label{b4:subsec:closed-external-internal}

内部 Hom 在范畴中可以继续张量、取极限或作其他对象构造；外部态射集合负责记录从一个对象到另一个对象的映射。两者通过单位对象联系起来。

\begin{proposition}\label{b4:prop:closed-unit-hom}
设 \(\mathcal C\) 为局部小闭幺半范畴，\(X,Y\) 为对象。有自然双射与自然同构
\[
\operatorname{Hom}_{\mathcal C}(X,Y)
\cong\operatorname{Hom}_{\mathcal C}(I,[X,Y]),
\qquad [I,Y]\cong Y.
\]
\end{proposition}
\begin{proof}
第一式由伴随取 \(A=I\)，再用 \(\lambda_X:I\otimes X\cong X\) 得到。第二式对任意 \(A\) 给出自然双射
\[
\operatorname{Hom}_{\mathcal C}(A,[I,Y])
\cong\operatorname{Hom}_{\mathcal C}(A\otimes I,Y)
\cong\operatorname{Hom}_{\mathcal C}(A,Y).
\]
最后一个双射由 \(\rho_A\) 给出且关于 \(A,Y\) 自然。可表对象的唯一性因而给出 \([I,Y]\cong Y\)，并由目标变量的自然性得到该同构关于 \(Y\) 自然。
\end{proof}

在集合中，从单点集到 \(Y^X\) 的映射就是选择一个函数。在模中，从 \(R\) 到 \(\operatorname{Hom}_R(X,Y)\) 的同态由 \(1\) 的像决定，仍对应一个 \(X\to Y\) 的模同态。单位对象使这个解释成立，而不是把内部 Hom 与态射集合当成同一种对象。

% !TeX root = ../../Book4.tex
% 纲要标题：### 5.3 求值、余求值与可对偶对象
% 纲要标签：b4:sec:monoidal-duals
% 纲要细目（写作范围）：
% #### 5.3.1 左右对偶与蛇形恒等式
% #### 5.3.2 对偶产生的伴随
% #### 5.3.3 可对偶模与有限生成投射性

\section{求值、余求值与可对偶对象}
\label{b4:sec:monoidal-duals}

有限维向量空间的恒等映射可以写成 \(v\mapsto\sum_i e_i^*(v)e_i\)。这个公式同时用了向量与泛函：先把有限个配对 \(e_i\otimes e_i^*\) 放在一起，再对输入向量求值，就恢复原向量。求值映射本身对无限维空间也存在；真正依赖有限性的，是能否用一个张量记录足够多的配对并恢复恒等映射。

\subsection{左右对偶与蛇形恒等式}
\label{b4:subsec:monoidal-left-right-duals}

\begin{definition}\label{b4:def:monoidal-dual}
设 \(\mathcal C\) 为幺半范畴，单位为 \(I\)，\(X,D\) 为其对象。给定态射
\[
e:D\otimes X\longrightarrow I,
\qquad b:I\longrightarrow X\otimes D.
\]
若下列两条复合都是恒等态射，则称 \((D,e,b)\) 为 \(X\) 的\term[zuo-duiou]{左对偶}[left dual]，常记 \(D=X^*\)：
\begin{equation}\label{b4:eq:monoidal-left-snakes}
\begin{aligned}
\rho_X\circ(1_X\otimes e)\circ\alpha_{X,D,X}
\circ(b\otimes1_X)\circ\lambda_X^{-1}&=1_X,\\
\lambda_D\circ(e\otimes1_D)\circ\alpha_{D,X,D}^{-1}
\circ(1_D\otimes b)\circ\rho_D^{-1}&=1_D.
\end{aligned}
\end{equation}
态射 \(e,b\) 分别称为求值与\term[yuqiuzhi]{余求值}[coevaluation]，上述条件称为蛇形恒等式。

若给定 \(e':X\otimes E\to I\)、\(b':I\to E\otimes X\)，且复合
\[
\begin{gathered}
\lambda_X\circ(e'\otimes1_X)\circ\alpha_{X,E,X}^{-1}
\circ(1_X\otimes b')\circ\rho_X^{-1},\\
\rho_E\circ(1_E\otimes e')\circ\alpha_{E,X,E}
\circ(b'\otimes1_E)\circ\lambda_E^{-1}
\end{gathered}
\]
分别为 \(1_X,1_E\)，则称 \((E,e',b')\) 为 \(X\) 的\term[you-duiou]{右对偶}[right dual]。若 \(X\) 同时有左对偶和右对偶，就称 \(X\) 为可对偶对象；若每个对象都可对偶，就称 \(\mathcal C\) 为\term[gangxing-fanchou]{刚性幺半范畴}[rigid monoidal category]。
\end{definition}

先暂略结合与单位约束，左对偶的两条蛇形复合分别是
\[
\begin{aligned}
X&\xrightarrow{\ b\otimes1_X\ }X\otimes D\otimes X
\xrightarrow{\ 1_X\otimes e\ }X,\\
D&\xrightarrow{\ 1_D\otimes b\ }D\otimes X\otimes D
\xrightarrow{\ e\otimes1_D\ }D.
\end{aligned}
\]
第一行的原输入在最右边，\(b\) 新产生左边的 \(X,D\)；\(e\) 消去中间的 \(D\) 与原输入，留下新产生的 \(X\)。第二行的原输入在最左边，\(b\) 新产生右边的 \(X,D\)；\(e\) 消去原输入与中间的 \(X\)，留下新产生的 \(D\)。要求这两种配对都恢复原输入，才得到恒等态射。在有限维向量空间中，这两行逐步成为
\[
\begin{aligned}
x&\longmapsto\sum_i e_i\otimes e_i^*\otimes x
\longmapsto\sum_i e_i^*(x)e_i=x,\\
\varphi&\longmapsto\sum_i\varphi\otimes e_i\otimes e_i^*
\longmapsto\sum_i\varphi(e_i)e_i^*=\varphi.
\end{aligned}
\]
定义中的 \(\lambda^{-1},\rho^{-1}\) 插入单位，\(\alpha\) 或 \(\alpha^{-1}\) 把将要求值的两个因子括在一起，最后的单位约束删去求值所得的 \(I\)。它们把上述配对过程写成定义域、陪域逐步相接的复合。

\ProseParagraphBreak

右对偶的两行则略写为
\[
X\longrightarrow X\otimes E\otimes X\longrightarrow X,
\qquad E\longrightarrow E\otimes X\otimes E\longrightarrow E.
\]
在第一行中，原输入 \(X\) 位于最左边，\(b'\) 在其右边产生 \(E,X\)，随后 \(e'\) 消去左边的 \(X,E\)；第二行中，原输入 \(E\) 位于最右边，\(b'\) 在其左边产生 \(E,X\)，随后消去右边的 \(X,E\)。因而两种侧别改变的是产生和求值的位置，不只是对偶对象的字母。

\ProseParagraphBreak

两个方向不能在一般幺半范畴中混用。左对偶的求值把 \(D\) 放在 \(X\) 左边，余求值却把它放在右边。\(D\) 是 \(X\) 的左对偶时，同一组态射使 \(X\) 成为 \(D\) 的右对偶；这还没有说明 \(X\) 自己有右对偶。在对称幺半范畴中，交换因子可以把两种方向互换。本节的侧别约定与\cite[Definitions 2.10.1--2.10.2, p. 40]{EtingofEtAl2015TensorCategories}一致。

\subsection{对偶产生的伴随}
\label{b4:subsec:monoidal-dual-adjunction}

余求值把一个空出的因子放进张量，求值再把它消去。这两个操作能将映射从张量的一侧转到另一侧。

\begin{proposition}\label{b4:prop:monoidal-dual-adjunction}
设 \(\mathcal C\) 为局部小幺半范畴，\((D,e,b)\) 为对象 \(X\) 的左对偶。对任意对象 \(A,B\)，有自然双射
\begin{equation}\label{b4:eq:monoidal-dual-adjunction}
\operatorname{Hom}_{\mathcal C}(A\otimes X,B)
\cong\operatorname{Hom}_{\mathcal C}(A,B\otimes D).
\end{equation}
因此 \(-\otimes D\) 是 \(-\otimes X\) 的右伴随。若 \(\mathcal C\) 还闭，则有自然同构 \([X,B]\cong B\otimes D\)。
\end{proposition}
\begin{proof}
双射及候选逆分别定义为
\[
\begin{aligned}
\Phi(f)&=(f\otimes1_D)\circ\alpha_{A,X,D}^{-1}
\circ(1_A\otimes b)\circ\rho_A^{-1},\\
\Psi(g)&=\rho_B\circ(1_B\otimes e)\circ\alpha_{B,D,X}
\circ(g\otimes1_X).
\end{aligned}
\]
这两式的定义域和陪域分别为 \(A\to B\otimes D\)、\(A\otimes X\to B\)。记\eqref{b4:eq:monoidal-left-snakes}的两条复合为 \(s_X,s_D\)。将 \(\Phi(f)\) 代入 \(\Psi\)，先用结合约束的自然性把 \(f\) 移到求值之后，再用五边形把四个因子的括号统一，三角等式删去单位约束，得到
\[
\Psi\Phi(f)=f\circ(1_A\otimes s_X)=f.
\]
这里剩下的三个因子按次序为 \(X,D,X\)，插入的是 \(b\)，消去的是 \(e\)，所以所用的是第一条蛇形恒等式。另一个复合按同样的自然性与五边形整理，剩下 \(D,X,D\)，因而
\[
\Phi\Psi(g)=(1_B\otimes s_D)\circ g=g.
\]
两式均只用了定义中的相容等式，证明互逆。

对 \(a:A'\to A\)、\(v:B\to B'\)，双函子性及约束自然性直接给出
\[
\Phi\bigl(v\circ f\circ(a\otimes1_X)\bigr)
=(v\otimes1_D)\circ\Phi(f)\circ a.
\]
这就是双射关于 \(A,B\) 的自然性。若还给内部 Hom，则 \([X,B]\) 与 \(B\otimes D\) 表示同一个函子 \(A\mapsto\operatorname{Hom}_{\mathcal C}(A\otimes X,B)\)，由可表对象唯一性得到所述自然同构。
\end{proof}

\begin{corollary}\label{b4:cor:monoidal-dual-unique}
设 \(\mathcal C\) 为局部小幺半范畴，\((D,e,b)\)、\((D',e',b')\) 都是 \(X\) 的左对偶。存在唯一同构 \(t:D\to D'\)，满足
\[
e'\circ(t\otimes1_X)=e,
\qquad (1_X\otimes t)\circ b=b'.
\]
右对偶也满足相应的唯一性。
\end{corollary}
\begin{proof}
在\cref{b4:prop:monoidal-dual-adjunction}中取 \(B=I\)，再用单位同构，得到关于 \(A\) 自然的双射
\[
\operatorname{Hom}_{\mathcal C}(A,D)
\cong\operatorname{Hom}_{\mathcal C}(A\otimes X,I).
\]
所以 \(D,D'\) 连同求值表示同一个函子，存在保持求值的唯一同构 \(t\)。把 \((1_X\otimes t)\circ b\) 代入 \(D'\) 所给的伴随逆，其像由第一条蛇形恒等式为 \(\lambda_X:I\otimes X\to X\)，与 \(b'\) 的像相同。伴随双射的单射性给出第二个等式。右对偶的证明使用相反的张量次序，完全得到同样的可表性和唯一性。
\end{proof}

这也说明，在闭范畴中存在 \([X,I]\) 还不够保证 \(X\) 可对偶。可对偶要求更强的事实：所有 \([X,B]\) 都由同一个对偶对象张量 \(B\) 得到。

\subsection{可对偶模与有限生成投射性}
\label{b4:subsec:monoidal-projective-duals}

在模范畴中，一项张量总是有限个纯张量的和。因此余求值若存在，会自动给出有限个生成元与有限个坐标泛函。

\begin{theorem}[可对偶模的判据]\label{b4:thm:dualizable-projective-modules}
设 \(R\) 为交换含幺环，\(P\) 为左 \(R\) 模。在对称幺半范畴 \((R\text{-}\mathbf{Mod},\otimes_R,R)\) 中，\(P\) 有左对偶，当且仅当它有限生成且投射；这也等价于 \(P\) 可对偶。此时可取对偶
\[
P^\vee=\operatorname{Hom}_R(P,R),
\qquad e(\varphi\otimes p)=\varphi(p).
\]
若 \(p_1,\ldots,p_r\in P\)、\(\varphi_1,\ldots,\varphi_r\in P^\vee\) 满足
\begin{equation}\label{b4:eq:monoidal-finite-dual-basis}
p=\sum_{i=1}^r\varphi_i(p)p_i\qquad(p\in P),
\end{equation}
则余求值由 \(b(1)=\sum_i p_i\otimes\varphi_i\) 给出。右对偶的求值与余求值交换相应因子的次序。
\end{theorem}
\begin{proof}
先设左对偶为 \((D,e,b)\)。将 \(b(1)\in P\otimes_RD\) 写成有限和 \(\sum_i p_i\otimes d_i\)，令 \(\varphi_i(p)=e(d_i\otimes p)\)。求值线性说明 \(\varphi_i\in P^\vee\)。第一条蛇形恒等式逐元素成为\eqref{b4:eq:monoidal-finite-dual-basis}，因此 \(p_i\) 生成 \(P\)。定义
\[
q:R^r\longrightarrow P,\quad(a_i)\longmapsto\sum_i a_ip_i,
\qquad j:P\longrightarrow R^r,\quad p\longmapsto(\varphi_i(p))_i.
\]
有 \(qj=1_P\)。于是 \(P\) 是有限秩自由模的直和因子，故投射。这里也可直接验证提升性质：给满射 \(M\to N\) 及 \(P\to N\)，先与 \(q\) 复合并提升 \(R^r\) 的各基向量，再与 \(j\) 复合，就得到原映射的提升。

第二条蛇形恒等式还给出 \(d=\sum_i e(d\otimes p_i)d_i\)。所以映射
\[
\begin{aligned}
D&\longrightarrow P^\vee,&d&\longmapsto(p\longmapsto e(d\otimes p)),\\
P^\vee&\longrightarrow D,&\varphi&\longmapsto\sum_i\varphi(p_i)d_i
\end{aligned}
\]
互逆：一个方向用第二条蛇形恒等式，另一个方向在 \(p\) 上取值后用\eqref{b4:eq:monoidal-finite-dual-basis}。这识别了对偶及其求值。

反过来，设 \(P\) 有限生成且投射。选有限秩自由满射 \(q:R^r\to P\)。投射性给出截面 \(j:P\to R^r\)。令 \(p_i=q(e_i)\)，\(\varphi_i\) 为 \(j\) 的第 \(i\) 坐标，便有\eqref{b4:eq:monoidal-finite-dual-basis}。上述求值保持平衡关系，余求值由其在 \(1\) 上的值唯一给出。第一条蛇形恒等式就是对偶基公式；第二条作用于 \(\varphi\in P^\vee\) 时成为
\[
\sum_i\varphi(p_i)\varphi_i=\varphi,
\]
因为两边在任意 \(p\) 上取值都等于 \(\varphi(p)\)。故得到左对偶。将求值、余求值中的两个因子互换，右对偶的两条蛇形恒等式也分别成为这两个已证公式，所以 \(P\) 可对偶。
\end{proof}

\begin{proposition}\label{b4:prop:projective-dual-internal-hom}
设 \(R\) 为交换含幺环，\(P\) 为有限生成投射 \(R\) 模，\(Y\) 为任意 \(R\) 模。自然映射
\[
\Theta_Y:Y\otimes_RP^\vee\longrightarrow\operatorname{Hom}_R(P,Y),
\qquad y\otimes\varphi\longmapsto(p\longmapsto\varphi(p)y)
\]
为同构。特别地，余求值的张量 \(b(1)\) 是 \(\Theta_P^{-1}(1_P)\)，不依赖对偶基的选择。
\end{proposition}
\begin{proof}
公式双线性且平衡。选满足\eqref{b4:eq:monoidal-finite-dual-basis}的一组数据，定义逆候选
\(f\mapsto\sum_i f(p_i)\otimes\varphi_i\)。先施行该候选再施行 \(\Theta_Y\)，在 \(p\) 上得到 \(\sum_i\varphi_i(p)f(p_i)=f(p)\)。另一个复合在 \(y\otimes\varphi\) 上得到
\[
\sum_i\varphi(p_i)y\otimes\varphi_i
=y\otimes\sum_i\varphi(p_i)\varphi_i=y\otimes\varphi.
\]
故两式互逆。\(\Theta_Y\) 本身未使用所选数据，所以其逆唯一，并且与 \(Y\) 中的同态相容。取 \(Y=P\) 即得最后的独立性结论。
\end{proof}

\begin{example}\label{b4:exm:dualizable-vector-spaces}
设 \(k\) 为域。向量空间可对偶，当且仅当有限维。对基 \(e_1,\ldots,e_r\) 及对偶基 \(e_1^*,\ldots,e_r^*\)，余求值为 \(1\mapsto\sum_i e_i\otimes e_i^*\)。无限维空间虽有代数对偶与求值，却没有满足蛇形恒等式的余求值：任何候选张量只涉及有限多个向量，第一条蛇形等式便迫使这些向量生成整个空间。因而全部向量空间的张量范畴闭但不刚性，有限维向量空间的张量范畴则刚性。
\end{example}

\begin{example}\label{b4:exm:torsion-no-dual}
在交换群范畴中，\(P=\mathbb Z/2\mathbb Z\) 不可对偶。事实上，每个 \(P\to\mathbb Z\) 的同态都为零：其像被 \(2\) 零化，而 \(\mathbb Z\) 无非零的这类元素。若存在左对偶 \(D\)，对固定 \(d\)，求值 \(p\mapsto e(d\otimes p)\) 因而为零，故整个求值为零。第一条蛇形复合只能为零，不可能为非零对象 \(P\) 的恒等映射。
\end{example}

% !TeX root = ../../Book4.tex
% 纲要标题：### 5.4 代数对象、模对象与富集
% 纲要标签：b4:sec:monoidal-algebras
% ---
% 纲要细目（写作范围）：
% #### 5.4.1 乘法与单位的对象形式
% #### 5.4.2 模对象与函子复合中的单子
% #### 5.4.3 内部复合与富集范畴

\section{代数对象、模对象与富集}
\label{b4:sec:monoidal-algebras}

普通代数的乘法是双线性映射，因而可以写成 \(A\otimes A\to A\)。其结合律比较三重乘积的两个求值过程，单位则来自基环到 \(A\) 的映射。这样写以后，运算的定义不再依赖元素记号：只要一个范畴有适当的张量与单位，就能讨论其中的代数与模。

\subsection{乘法与单位的对象形式}
\label{b4:subsec:monoidal-algebra-objects}

\begin{definition}\label{b4:def:monoidal-algebra-object}
设 \(\mathcal C\) 为幺半范畴，单位为 \(I\)。给定对象 \(A\) 及态射 \(m:A\otimes A\to A\)、\(u:I\to A\)。若满足
\begin{equation}\label{b4:eq:monoidal-algebra-laws}
\begin{aligned}
m\circ(m\otimes1_A)&=m\circ(1_A\otimes m)\circ\alpha_{A,A,A},\\
m\circ(u\otimes1_A)&=\lambda_A,\qquad
m\circ(1_A\otimes u)=\rho_A,
\end{aligned}
\end{equation}
则称 \((A,m,u)\) 为 \(\mathcal C\) 中的\term[daishu-duixiang]{代数对象}[algebra object]，也称幺半群对象。对两个代数对象，若态射 \(f:A\to B\) 满足
\[
f\circ m_A=m_B\circ(f\otimes f),\qquad f\circ u_A=u_B,
\]
则称 \(f\) 为代数对象的同态。在辫状幺半范畴中，若还满足 \(m\circ c_{A,A}=m\)，则称代数对象为交换的。
\end{definition}

第一式从 \((A\otimes A)\otimes A\) 出发：左边先乘前两个因子，再与第三个相乘；右边先用 \(\alpha\) 改括号，乘后两个因子，再与第一个相乘。它要求这两种乘法过程相同。第二行则先以 \(u\) 把单位对象变成代数中的单位，再与原输入相乘，所得映射分别等于 \(\lambda_A:I\otimes A\to A\)、\(\rho_A:A\otimes I\to A\)。由于定义域分别为 \(I\otimes A\)、\(A\otimes I\)，这里不能直接写成 \(1_A\)。这一形式可对照\cite[Definition 7.8.1, p. 141]{EtingofEtAl2015TensorCategories}；其定义所在的张量范畴带有其他有限性条件，但上述乘法公理本身只需幺半结构。

\begin{proposition}\label{b4:prop:monoidal-ordinary-algebras}
设 \(R\) 为交换含幺环。在 \((R\text{-}\mathbf{Mod},\otimes_R,R)\) 中，代数对象及其同态恰好给出含幺结合 \(R\) 代数及其保幺 \(R\) 代数同态。其交换代数对象恰好是交换 \(R\) 代数。在 \((\mathbf{Set},\times,\{*\})\) 中，代数对象及同态恰好是幺半群及幺半群同态。
\end{proposition}
\begin{proof}
给 \(R\)-模代数对象，定义 \(ab=m(a\otimes b)\)、\(1_A=u(1)\)。张量积泛性质使该乘法双线性。将\eqref{b4:eq:monoidal-algebra-laws}作用于 \((a\otimes b)\otimes c\)，得到 \((ab)c=a(bc)\)；作用于 \(1\otimes a\)、\(a\otimes1\)，得到 \(1_Aa=a=a1_A\)。又因 \(u\) 线性，\(u(r)=r1_A\)，于是
\[
u(r)a=ra=au(r),\qquad u(r)u(s)=u(rs).
\]
所以 \(u:R\to A\) 是取值于中心的保幺环同态，得到 \(R\) 代数。标量作用 \(ra\) 原先已由 \(R\) 模 \(A\) 给定；\(u(r)a=ra\) 证明沿 \(u\) 诱导的标量作用与这个原有作用相同。

反过来，含幺结合 \(R\) 代数的乘法双线性且平衡，故诱导唯一的 \(m:A\otimes_RA\to A\)。结构同态 \(u\) 线性；结合、单位公理在纯张量上成立，故给出全部所需态射等式。两构造在元素与结构映射上互逆。代数对象同态的两个条件正是保持乘法与单位，且其本身是 \(R\) 线性映射，所以恰是保幺 \(R\) 代数同态。交换约束把 \(a\otimes b\) 换成 \(b\otimes a\)，故附加条件等价于 \(ab=ba\)。

对集合，\(m\) 是二元运算，\(u\) 指定元素 \(u(*)\)。同样的等式直接成为结合律与左右单位律，同态条件也直接成为幺半群同态的条件。
\end{proof}

交换性取决于所选辫状结构。例如，在分次向量空间的符号对称结构中，对齐次元素 \(a\in A_m\)、\(b\in A_n\)，条件 \(m\circ c=m\) 写成
\[
ab=(-1)^{mn}ba.
\]
它允许两个奇次数元素的乘法带负号，因而与不分次数的普通交换性不同。对奇次数元素 \(a\)，直接得到的是 \(a^2=-a^2\)，即 \(2a^2=0\)。若 \(2\) 可逆，则由此推出 \(a^2=0\)；特征 \(2\) 时符号为正，不能这样推出平方为零。

\subsection{模对象与函子复合中的单子}
\label{b4:subsec:monoidal-module-objects}

一个代数对象要作用在另一个对象上，只须把乘法的第二个因子改成该对象，并保留同样的结合与单位要求。

\begin{definition}\label{b4:def:monoidal-module-object}
设 \((A,m,u)\) 为幺半范畴 \(\mathcal C\) 中的代数对象。给对象 \(M\) 及态射 \(a:A\otimes M\to M\)。若满足
\[
a\circ(m\otimes1_M)=a\circ(1_A\otimes a)\circ\alpha_{A,A,M},
\qquad a\circ(u\otimes1_M)=\lambda_M,
\]
则称 \((M,a)\) 为左 \(A\)-\term[mo-duixiang]{模对象}[module object]。若 \(f:M\to N\) 满足 \(f\circ a_M=a_N\circ(1_A\otimes f)\)，则称其为模对象同态。右模对象以 \(M\otimes A\to M\) 的作用定义，其结合等式为
\(a\circ(a\otimes1_A)=a\circ(1_M\otimes m)\circ\alpha_{M,A,A}\)，单位等式为 \(a\circ(1_M\otimes u)=\rho_M\)。
\end{definition}

从 \((A\otimes A)\otimes M\) 出发，作用结合式的左边先乘两个 \(A\) 因子，再作用于 \(M\)；右边改括号后，先让第二个 \(A\) 因子作用于 \(M\)，再让第一个作用。单位式要求插入代数单位后的作用恢复原输入。在交换环 \(R\) 的模范畴中，左模对象恢复通常的左 \(A\) 模。确切地，作用态射对应 \(R\) 双线性作用 \((x,v)\mapsto xv\)，两条等式成为 \((xy)v=x(yv)\)、\(1_Av=v\)。反过来，普通左 \(A\) 模沿 \(R\to A\) 取得 \(R\) 模结构，其作用双线性且平衡，故诱导模对象结构。两种同态条件也都成为 \(f(xv)=xf(v)\)。在集合的积结构中，这一过程给出幺半群在集合上的作用，而不是预先带加法的模。

\begin{proposition}\label{b4:prop:monads-as-monoidal-algebras}
设 \(\mathcal C\) 为小范畴。自函子范畴 \(\operatorname{Fun}(\mathcal C,\mathcal C)\) 以函子复合作为张量、以恒等函子作为单位，是严格幺半范畴。其代数对象恰好是 \(\mathcal C\) 上的单子。
\end{proposition}
\begin{proof}
对象的张量为 \(F\otimes G=F\circ G\)，自然变换的张量为水平复合。水平复合保持垂直复合与恒等自然变换，故确实给出双函子；这是自然变换交换律的内容。函子复合及自然变换的水平复合严格结合，恒等函子严格为单位，所以约束均可取恒等。

代数对象的对象现在是自函子 \(T\)。乘法 \(m:T\otimes T\to T\) 变成自然变换 \(\mu:T\circ T\Rightarrow T\)，单位 \(u:I\to T\) 变成 \(\eta:1_{\mathcal C}\Rightarrow T\)。左右张量 \(\mu\) 所得的分量分别是 \(\mu_{TX}:T^3X\to T^2X\)、\(T(\mu_X):T^3X\to T^2X\)，张量 \(\eta\) 所得的分量则是 \(\eta_{TX}\)、\(T(\eta_X)\)。因此代数对象公理逐分量写成
\[
\mu_X\circ\mu_{TX}=\mu_X\circ T(\mu_X),
\qquad
\mu_X\circ\eta_{TX}=1_{TX}=\mu_X\circ T(\eta_X),
\]
这恰好是\cref{b4:def:monad}的单子公理。反过来，每个单子给出这些相同的数据与等式，故两种结构完全对应。
\end{proof}

这一个例子中的张量是函子复合，通常不能交换两个因子。幺半结构能够容纳这种情形，因此代数对象不以辫状或对称结构为前提。

\subsection{内部复合与富集范畴}
\label{b4:subsec:monoidal-enrichment}

在线性代数中，同态不仅组成集合，还能相加；复合是双线性运算。因此把 Hom 作为对象，能够同时保留态射和态射之间的线性结构。内部 Hom 给出一种实现这一想法的方式。要把普通复合 \((g,f)\mapsto g\circ f\) 写成内部 Hom 之间的态射，可以先在输入上连续求值，再用转置确定所需态射；在模的例子中，这个求值过程就是 \((g,f,x)\mapsto g(f(x))\)。

\begin{proposition}\label{b4:prop:closed-internal-composition}
设 \(\mathcal C\) 为局部小闭幺半范畴。求值映射唯一确定内部复合
\[
\gamma_{X,Y,Z}:[Y,Z]\otimes[X,Y]\longrightarrow[X,Z]
\]
及单位 \(j_X:I\to[X,X]\)。其中 \(\gamma\) 是从 \(([Y,Z]\otimes[X,Y])\otimes X\) 到 \(Z\) 的下列态射的转置：
\[
\operatorname{ev}_{Y,Z}\circ
(1_{[Y,Z]}\otimes\operatorname{ev}_{X,Y})\circ
\alpha_{[Y,Z],[X,Y],X},
\]
而 \(j_X\) 是 \(\lambda_X:I\otimes X\to X\) 的转置。这些态射满足结合与左右单位等式。因此 \([X,X]\) 带有典范的代数对象结构。
\end{proposition}
\begin{proof}
转置的存在唯一性先给出 \(\gamma,j_X\)。对四个对象 \(W,X,Y,Z\)，比较
\[
\begin{gathered}
\gamma_{W,X,Z}\circ(\gamma_{X,Y,Z}\otimes1_{[W,X]}),\\
\gamma_{W,Y,Z}\circ(1_{[Y,Z]}\otimes\gamma_{W,X,Y})
\circ\alpha_{[Y,Z],[X,Y],[W,X]}.
\end{gathered}
\]
两者的定义域都是 \(([Y,Z]\otimes[X,Y])\otimes[W,X]\)。在模的例子中，给定 \(f:W\to X\)、\(g:X\to Y\)、\(h:Y\to Z\) 与 \(w\in W\)，两种复合都取值为 \(h(g(f(w)))\)。一般情形同样用连续求值比较：将这两个态射与 \(W\) 张量后求值，第一种复合先在 \(W\) 上求值到 \(X\)，再到 \(Y\)，最后到 \(Z\)；第二种也依次使用同样的三个求值。结合约束的五边形保证两种重括号到
\([Y,Z]\otimes([X,Y]\otimes([W,X]\otimes W))\) 相同，故所得态射相等。转置的唯一性因而给出结合等式。

将 \(j_Y\) 与 \([X,Y]\) 作内部复合，其在 \(X\) 上的求值仍是 \(\operatorname{ev}_{X,Y}\)，因为 \(j_Y\) 的转置是单位映射 \(I\otimes Y\to Y\)。三角等式删除这项单位，转置唯一性给出左单位律。将 \(j_X\) 放在另一侧，同样在 \(X\) 上求值得到原求值，给出右单位律。最后取三个对象都为 \(X\)，这些等式就是代数对象公理。
\end{proof}

在交换环 \(R\) 的模范畴中，内部复合为 \(g\otimes f\mapsto g\circ f\)，单位把 \(1\) 送到恒等同态。因此 \([P,P]=\operatorname{End}_R(P)\) 的代数对象结构就是通常的自同态代数。这也解释乘法次序为何按先 \(f\)、后 \(g\) 排列。
\symidx{\(\operatorname{End}_R(P)\)：模的自同态环}

普通局部小范畴为每对对象给出一个 Hom 集合。若把它换成 \(\mathcal V\) 中的对象，复合所用的集合笛卡尔积便换成 \(\mathcal V\) 的张量积，恒等元素也换成从单位对象 \(I\) 出发的态射。下面的定义保留结合与单位要求，以这样的 Hom 对象作为数据。

\begin{definition}\label{b4:def:enriched-category}
设 \(\mathcal V\) 为局部小幺半范畴，单位为 \(I\)。给定一个对象类，对每对对象 \(X,Y\) 给 \(\mathcal V\) 中的对象 \(\mathcal H(X,Y)\)，以及态射
\[
\gamma_{X,Y,Z}:\mathcal H(Y,Z)\otimes\mathcal H(X,Y)\longrightarrow\mathcal H(X,Z),
\qquad j_X:I\longrightarrow\mathcal H(X,X).
\]
若对任意 \(W,X,Y,Z\)，满足
\[
\begin{aligned}
\gamma_{W,X,Z}\circ(\gamma_{X,Y,Z}\otimes1)
={}&\gamma_{W,Y,Z}\circ(1\otimes\gamma_{W,X,Y})
\circ\alpha,\\
\gamma_{X,Y,Y}\circ(j_Y\otimes1)&=\lambda_{\mathcal H(X,Y)},\\
\gamma_{X,X,Y}\circ(1\otimes j_X)&=\rho_{\mathcal H(X,Y)},
\end{aligned}
\]
其中第一式的 \(\alpha\) 是 \(\mathcal H(Y,Z),\mathcal H(X,Y),\mathcal H(W,X)\) 的结合约束，恒等态射各在相应 Hom 对象上，则称这些数据为一个\term[fujifanchou]{富集范畴}[enriched category]，并称它富集于 \(\mathcal V\)。这里 Hom 对象不必是某个既有范畴的内部 Hom，而是所给的新数据。
\end{definition}
\symidx{\(\mathcal H(X,Y)\)：富集范畴中的态射对象}

当 \(\mathcal V=\mathbf{Ab}\)，张量为 \(\otimes_{\mathbb Z}\) 时，内部复合就是群同态 \(\mathcal H(Y,Z)\otimes\mathcal H(X,Y)\to\mathcal H(X,Z)\)，等价于双可加复合；\(\mathbb Z\to\mathcal H(X,X)\) 由恒等态射决定。因此这恰好恢复第四章预加性范畴的 Hom 群及复合。改用 \(k\)-向量空间，则恢复 \(k\) 线性范畴：复合对两个变量分别 \(k\) 线性。两者再要求零对象及有限双积存在，才得到加性范畴。

\ProseParagraphBreak

一般富集数据也给出普通范畴：取从 \(X\) 到 \(Y\) 的态射集合为 \(\operatorname{Hom}_{\mathcal V}(I,\mathcal H(X,Y))\)，即用从单位对象出发的态射来取 Hom 对象的“元素”。对 \(f:I\to\mathcal H(X,Y)\)、\(g:I\to\mathcal H(Y,Z)\)，张量 \(g\otimes f\) 的源是 \(I\otimes I\)，需要先接上单位同构的逆 \(\delta=\lambda_I^{-1}:I\to I\otimes I\)。在 Ab 或向量空间中，\(\delta\) 就是 \(1\mapsto1\otimes1\)。为与 \(\mathcal V\) 本身的态射复合区别，验证期间把新复合记为 \(\star\)：
\[
 g\star f\coloneqq\gamma_{X,Y,Z}\circ(g\otimes f)\circ\delta
 :I\longrightarrow\mathcal H(X,Z),
\]
恒等态射取为 \(j_X:I\to\mathcal H(X,X)\)。

\ProseParagraphBreak

若幺半结构严格，则 \(I\otimes I=I\)、\(\delta=1_I\)，新复合的结合律直接来自富集结合律。一般情形还需核对插入单位和改变括号的复合相容。单位约束满足 \(\lambda_I=\rho_I\)\cite[Corollary 2.2.5, p. 24]{EtingofEtAl2015TensorCategories}。将三角等式取 \(X=Y=I\) 并取逆，再前复合 \(\delta\)，得到
\[
 \alpha_{I,I,I}\circ(\delta\otimes1_I)\circ\delta
 =(1_I\otimes\delta)\circ\delta.
\]
结合 \(\alpha\) 的自然性与富集结合等式，对可相接的 \(f,g,h\) 得到 \((h\star g)\star f=h\star(g\star f)\)：两边插入三个单位时的差别恰由上式消去。若简记 \(H=\mathcal H(X,Y)\)，富集单位等式以及 \(\lambda,\rho\) 的自然性给出
\[
\begin{aligned}
 j_Y\star f&=\lambda_H\circ(1_I\otimes f)\circ\delta
 =f\circ\lambda_I\circ\delta=f,\\
 f\star j_X&=\rho_H\circ(f\otimes1_I)\circ\delta
 =f\circ\rho_I\circ\delta=f.
\end{aligned}
\]
这就验证了普通范畴的结合律与左右单位律；以后可把 \(\star\) 仍记为通常的复合符号 \(\circ\)。

\ProseParagraphBreak

在 Ab 和向量空间例子中，这正是取回 Hom 对象的元素。\chapref{b4:ch:20}的 DG 范畴把 Hom 对象改为复形，复合还必须与微分相容。

\ProseParagraphBreak
若 Hom 对象取为小范畴，张量取笛卡尔积，严格的复合与单位函子便给出严格2范畴。继续让 Hom 取值于前一层的严格范畴，就得到\secref{b4:sec:F03}的递归定义。


\begin{chaptersummary}
幺半结构不只规定一种对象运算，还指定结合与单位同构，并要求不同操作路径相容。辫状结构进一步规定因子的交换方式；对称结构要求交换两次恢复原映射。集合的积与交换环模的张量都闭，但它们的单位、内部 Hom 以及可对偶对象并不相同。

\ProseParagraphBreak

可对偶性用求值和余求值恢复恒等态射，因此使内部 Hom 化为与固定对偶对象的张量。交换环模中的余求值只能使用有限个纯张量，这一有限性恰好产生有限生成投射模的对偶基。代数对象把乘法与单位写成相容态射，模对象将同一条件用于作用；富集则让 Hom 本身保留群、向量空间或其他对象结构，复合也在这一结构中进行。
\end{chaptersummary}

\BookExercises

\begin{exercise}\label{b4:ex:monoidal-extension-scalars}
设 \(R\to S\) 为交换含幺环同态。证明标量扩张函子 \(F(M)=S\otimes_RM\) 带有强幺半结构，其中比较映射为
\[
\begin{gathered}
(S\otimes_RM)\otimes_S(S\otimes_RN)\longrightarrow S\otimes_R(M\otimes_RN),
\\
(s\otimes m)\otimes(t\otimes n)\longmapsto st\otimes(m\otimes n).
\end{gathered}
\]
写出逆映射与单位比较，验证它们保持结合、单位与交换约束。
\end{exercise}
\begin{solution}
两个内层的 \(R\) 平衡关系都由结构同态保持；外层将 \(S\) 系数从一个因子移到另一个因子的平衡关系，则由 \(st\) 不变保持。因此比较映射良定义并为 \(S\) 线性。逆为
\[
s\otimes(m\otimes n)\longmapsto(s\otimes m)\otimes(1\otimes n).
\]
把 \(rm\otimes n\) 换成 \(m\otimes rn\)，其像由内层及外层平衡关系相等，故逆也良定义。一个复合直接恢复 \(s\otimes(m\otimes n)\)；另一个复合将 \((st\otimes m)\otimes(1\otimes n)\) 化为 \((s\otimes m)\otimes(t\otimes n)\)。单位比较为 \(S\to S\otimes_RR\)、\(s\mapsto s\otimes1\)，逆为 \(s\otimes r\mapsto s r\)，其中 \(r\) 通过结构同态作用。

结合的两条路径在三个纯张量上都把三个 \(S\) 系数相乘，并按同样次序留下 \(m,n,p\)。单位的两条路径都使单位因子系数作用在另一因子上。交换的两条路径分别留下 \(n\otimes m\) 并乘 \(st=ts\)。因此各相容等式在纯张量上成立；纯张量生成相应模，故得到所需强幺半结构，且它尊重对称结构。
\end{solution}

\begin{exercise}\label{b4:ex:direct-sum-not-closed}
设 \(R\) 为环，固定左 \(R\) 模 \(M\)。证明函子 \(-\oplus M\) 有右伴随，当且仅当 \(M=0\)。据此判断以直接和为幺半运算的模范畴是否闭。
\end{exercise}
\begin{solution}
若有右伴随 \(G\)，则对任意 \(N\)，取伴随的源为零模，得到
\[
\operatorname{Hom}_R(M,N)
\cong\operatorname{Hom}_R(0,G(N)).
\]
右边只有零同态，故左边也只能有一个同态。取 \(N=M\) 得到 \(1_M=0\)，因此 \(M=0\)。反过来，\(-\oplus0\) 自然同构于恒等函子，以恒等函子作右伴随即可。只要模范畴存在非零对象，其直接和幺半结构就不闭；对零环的只有零模的范畴则闭。
\end{solution}

\begin{exercise}\label{b4:ex:cartesian-dualizable-sets}
证明在集合范畴的笛卡尔积幺半结构中，一个集合有左对偶，当且仅当它是单点集。空集是否例外？
\end{exercise}
\begin{solution}
若 \(D\) 是 \(X\) 的左对偶，余求值 \(\{*\}\to X\times D\) 指定一对 \((x_0,d_0)\)，所以两集合非空。求值 \(D\times X\to\{*\}\) 唯一。第一条蛇形复合把每个 \(x\) 都送到 \(x_0\)，而它须为恒等，故 \(X=\{x_0\}\)。第二条同样迫使 \(D=\{d_0\}\)。反过来，两个单点集之间的唯一映射给出求值和余求值，两个蛇形复合均为恒等。空集没有余求值，因为其与任意集合的积为空，所以不是例外。
\end{solution}

\begin{exercise}\label{b4:ex:dual-of-tensor-modules}
设 \(R\) 为交换含幺环，\(P,Q\) 为有限生成投射 \(R\) 模。证明 \(P\otimes_RQ\) 可对偶，并将其对偶识别为 \(Q^\vee\otimes_RP^\vee\)。写出求值及余求值，不仅给出存在性。
\end{exercise}
\begin{solution}
选对偶基 \((p_i,\varphi_i)\)、\((q_j,\psi_j)\)。求值定义为
\[
(\psi\otimes\varphi)\otimes(p\otimes q)\longmapsto\varphi(p)\psi(q),
\]
余求值把 \(1\) 送到 \(\sum_{i,j}(p_i\otimes q_j)\otimes(\psi_j\otimes\varphi_i)\)。各式 \(R\) 多线性且保持平衡关系，故良定义。第一个蛇形复合在 \(p\otimes q\) 上为
\[
\sum_{i,j}\varphi_i(p)\psi_j(q)(p_i\otimes q_j)
=\left(\sum_i\varphi_i(p)p_i\right)\otimes
\left(\sum_j\psi_j(q)q_j\right)=p\otimes q.
\]
第二个在 \(\psi\otimes\varphi\) 上成为
\(\sum_{i,j}\varphi(p_i)\psi(q_j)\psi_j\otimes\varphi_i=\psi\otimes\varphi\)。故给出左对偶，交换因子也给出右对偶。由对偶唯一性及\cref{b4:thm:dualizable-projective-modules}，它识别为 \(\operatorname{Hom}_R(P\otimes_RQ,R)\)，并同时证明张量积有限生成投射。
\end{solution}

\begin{exercise}\label{b4:ex:dual-morphism-compatibility}
设 \(R\) 为交换含幺环，\(P,Q\) 为有限生成投射 \(R\) 模，\(f:P\to Q\)。令 \(f^\vee:Q^\vee\to P^\vee\)、\(\psi\mapsto\psi\circ f\)。证明对偶反转复合，并验证余求值的相容公式
\[
(f\otimes1_{P^\vee})\circ b_P
=(1_Q\otimes f^\vee)\circ b_Q:R\longrightarrow Q\otimes_RP^\vee.
\]
\end{exercise}
\begin{solution}
对 \(g:Q\to T\) 及 \(\chi\in T^\vee\)，有 \((gf)^\vee(\chi)=\chi\circ g\circ f=f^\vee g^\vee(\chi)\)，而恒等的对偶仍为恒等，所以得到反变函子。相容公式左边在 \(1\) 处的张量为 \(\sum_i f(p_i)\otimes\varphi_i\)，经\cref{b4:prop:projective-dual-internal-hom}的同构映为 \(f\)。右边为 \(\sum_jq_j\otimes(\psi_j\circ f)\)，其对应同态在 \(p\) 上的值为 \(\sum_j\psi_j(f(p))q_j=f(p)\)，也对应 \(f\)。同构的单射性给出两张量相等，故整个映射相等。
\end{solution}

\begin{exercise}\label{b4:ex:dualizable-trace-cyclicity}
设 \(R\) 为交换含幺环，\(P\) 为有限生成投射 \(R\) 模。对 \(h\in\operatorname{End}_R(P)\)，用求值、交换与余求值定义
\[
\operatorname{tr}_P(h)=
\bigl(e_P\circ c_{P,P^\vee}\circ(h\otimes1_{P^\vee})\circ b_P\bigr)(1)\in R.
\]
写出对偶基公式，并证明对有限生成投射 \(P,Q\) 及 \(f:P\to Q\)、\(g:Q\to P\)，有 \(\operatorname{tr}_P(gf)=\operatorname{tr}_Q(fg)\)。当 \(P=R^r\) 时与矩阵迹比较。
\end{exercise}
\begin{solution}
沿各映射计算得 \(\operatorname{tr}_P(h)=\sum_i\varphi_i(h(p_i))\)。余求值及其余态射均已确定，不依赖对偶基，故这一公式也不依赖选择。对 \(Q\) 选对偶基 \((q_j,\psi_j)\)，则
\[
\begin{aligned}
\operatorname{tr}_P(gf)
&=\sum_{i,j}\psi_j(f(p_i))\varphi_i(g(q_j))\\
&=\sum_j\psi_j\left(f\left(\sum_i\varphi_i(g(q_j))p_i\right)\right)
=\sum_j\psi_j(fg(q_j))=\operatorname{tr}_Q(fg).
\end{aligned}
\]
将标量顺序交换用了 \(R\) 交换。对 \(R^r\) 取标准基及坐标泛函，公式就是矩阵的对角元之和，恒等映射的迹为 \(r\cdot1_R\)。
\end{solution}

\begin{exercise}\label{b4:ex:graded-odd-squares}
设 \(k\) 为特征不为 \(2\) 的域，\(A\) 为符号对称结构下的 \(\mathbb Z\) 分次交换代数对象。证明每个奇次数齐次元素的平方为零。特征为 \(2\) 时构造反例。
\end{exercise}
\begin{solution}
若 \(a\) 的次数 \(m\) 为奇数，则 \(a^2=(-1)^{m^2}a^2=-a^2\)，故 \(2a^2=0\)。特征不为 \(2\) 时 \(2\) 可逆，因此 \(a^2=0\)。特征为 \(2\) 时取多项式代数 \(k[x]\)，赋予 \(\deg x=1\)。符号 \(-1=1\)，所以普通交换乘法满足分次交换条件，却有 \(x^2\ne0\)。
\end{solution}

\begin{exercise}\label{b4:ex:one-object-enrichment}
证明只有一个对象的 Ab 富集范畴等价于给定一个含幺环；只有一个对象的 \(k\)-向量空间富集范畴等价于给定一个含幺结合 \(k\) 代数。富集范畴的结合律是否迫使环交换？
\end{exercise}
\begin{solution}
设唯一对象为 \(*\)，其 Hom 对象为 \(H\)。在 Ab 中，复合 \(H\otimes_{\mathbb Z}H\to H\) 等价于双可加乘法，单位 \(\mathbb Z\to H\) 指定元素 \(1\)。富集的结合与单位公理就是环的结合与单位律，双可加性给出左右分配律。反过来，任意含幺环都给出这些数据，故互相恢复。在向量空间中将双可加换成 \(k\) 双线性、单位域换成 \(k\)，便得到含幺结合 \(k\) 代数。公理只要求结合，不要求交换；例如矩阵环 \(M_2(k)\) 给出只有一个对象的富集范畴，而 \(E_{12}E_{21}=E_{11}\ne E_{22}=E_{21}E_{12}\)。
\end{solution}

\begin{exercise}\label{b4:ex:direct-sum-algebra-objects}
设 \(R\) 为环。在以 \(\oplus\) 为幺半运算、零模为单位的左 \(R\) 模范畴中，分类代数对象。对给定代数对象 \(A\)，证明左模对象等价于给定一个左 \(R\) 模 \(M\) 和同态 \(t:A\to M\)。
\end{exercise}
\begin{solution}
单位 \(0\to A\) 唯一。乘法 \(m:A\oplus A\to A\) 由在两个直和分量上的限制决定，而左右单位公理要求两个限制都是 \(1_A\)。因此唯一可能的乘法为 \(m(a,b)=a+b\)，它满足结合律。每个 \(R\) 模恰好有这一种代数对象结构，每个模同态都保持它。

作用 \(a:A\oplus M\to M\) 的单位公理要求在 \(M\) 分量上为恒等，故它唯一写成 \(a(x,v)=t(x)+v\)。任意 \(R\) 线性 \(t\) 都满足结合公理，因为两边在 \((x,y,v)\) 上分别为 \(t(x+y)+v\)、\(t(x)+t(y)+v\)。于是给定左模对象等价于给定这样的 \(t\)，其同态 \(f:M\to N\) 的条件为 \(f t_M=t_N\)。这与张量幺半结构下的普通代数和普通模不同。
\end{solution}

\begin{exercise}\label{b4:ex:representation-duals}
设 \(G\) 为群，\(k\) 为域，\(V\) 为有限维 \(k\) 向量空间，给定表示 \(\rho:G\to\operatorname{GL}_k(V)\)。在 \(V^*\) 上定义 \((g\varphi)(v)=\varphi(g^{-1}v)\)。证明这给出表示，并且通常的求值与余求值对 \(G\) 等变，其中 \(k\) 取平凡作用、张量取对角作用。是否需要 \(G\) 有限或 \(\operatorname{char}k\) 不整除群阶？
\end{exercise}
\begin{solution}
有 \(1\varphi=\varphi\)，且
\((g(h\varphi))(v)=\varphi(h^{-1}g^{-1}v)=((gh)\varphi)(v)\)，故为表示。求值满足 \((g\varphi)(gv)=\varphi(v)\)，所以等变。余求值张量 \(\sum_i e_i\otimes e_i^*\) 在典范同构 \(V\otimes_kV^*\cong\operatorname{End}_k(V)\) 下对应 \(1_V\)。对角作用下，\(v\otimes\varphi\) 对应的算子变为 \(w\mapsto\varphi(g^{-1}w)gv\)，正是原算子的 \(\rho(g)\) 共轭。因此恒等算子固定，余求值张量也固定，给出等变映射 \(k\to V\otimes_kV^*\)。蛇形等式已在底向量空间中成立，也就在表示之间成立。整个论证只需要有限维及群元素可逆，不需要 \(G\) 有限，也不需要群阶在域中可逆。
\end{solution}
